\section{Advising the wealthy and transferring wealth}

\begin{qfmethod}[Citizenship, nationality and residency in private wealth planning: US citizenship-based taxation and exit tax, the US substantial presence test, the UK statutory residence test, the UK's residence-based regimes from 2025, Singapore's residence rule and treaty tie-breakers]\label{meth:residency-citizenship}
\qffield{Statement} \emph{Citizenship against residence.} The United States taxes its citizens wherever they live; most countries tax by residence. A US citizen abroad therefore stays in the US system unless they renounce, and the UK--US treaty preserves that right (Art.~1(4)) \citep[Art.~1(4)]{uk-us-double-taxation-convention-2001}.
\begin{itemize}
\item \emph{US residence for non-citizens}: a green card or the \emph{substantial presence test}, 31 days in the year and at least 183 days counting all days this year, a third of last year's and a sixth of the year before's, with a closer-connection exception for those present fewer than 183 days \citep[sec.~7701(b)(3)]{us-internal-revenue-code-26}.
\item \emph{US exit tax}: a \emph{covered expatriate} is one with average annual net income tax above an indexed threshold (USD~211,000 for 2026), net worth of 2 million or more, or who fails to certify five years of compliance \citep[sec.~877(a)(2)]{us-internal-revenue-code-26}; \citep[sec.~4.37]{irs-rev-proc-2025-32}. All their property is treated as sold the day before expatriation, with gains above an exclusion (USD~910,000 for 2026) taxed \citep[sec.~877A(a)]{us-internal-revenue-code-26}; \citep[sec.~4.38]{irs-rev-proc-2025-32}.
\item \emph{UK residence}: the statutory residence test: 183 days or more is residence; otherwise automatic overseas tests (fewer than 16 days if resident in one of the last three years, fewer than 46 days if not, or full-time work abroad), automatic UK tests (a UK home, full-time UK work) and, failing those, a sufficient-ties test whose required number of ties falls as days rise \citep[secs.~2--4, Tables A and B]{hmrc-rdr3-2026}.
\item \emph{UK regimes from 6 April 2025}: domicile no longer matters for income tax; new arrivals who were not UK resident in the previous 10 years can claim relief on foreign income and gains for their first four years of residence \citep[policy summary]{govuk-nondom-reform-2024}. Inheritance tax on foreign assets follows long-term residence: 10 of the previous 20 tax years, with the tail after leaving lasting 3 to 10 years \citep[IHTM47020]{hmrc-ihtm-2026}.
\item \emph{Singapore}: a foreigner is tax resident after at least 183 days in the previous calendar year, or three consecutive years, among other tests \citep[residency rules]{iras-tax-residency-2026}.
\item \emph{Dual residence}: treaties assign residence by permanent home, then centre of vital interests, then habitual abode, then nationality, then mutual agreement (UK--US treaty, Art.~4(4)).
\end{itemize}
Nationality also determines passports, visas and the right to live and work in a place, and can be acquired through residence, descent or investment programmes (Interpretation).

\qffield{Practice} \emph{Maya Chen} (original; \texttt{private-wealth-checks.py}, part~11), 45, US citizen, an executive in New York, is offered a five-year post in London and might later retire to Singapore.
\begin{itemize}
\item \emph{UK}: she has not been UK resident in the last 10 years, so in her first four UK tax years her foreign income and gains (for example, US investment income) can be relieved from UK tax if she claims; UK earnings are taxed. After 10 UK-resident years her worldwide estate would be in UK inheritance tax.
\item \emph{US}: she remains taxable as a citizen and claims credits for UK tax; renouncing would make her a covered expatriate (net worth above 2 million): with 3.5 million of unrealised gains, 2.59 million would be taxed, about 616,000 at 23.8 percent.
\item \emph{Days}: on leaving London, spending 100 days in the UK in a year would keep her UK resident if she kept two ties (Table~A); a returning US non-citizen spouse with 150, 120 and 120 days in successive years would meet the US test (210 weighted days).
\item \emph{Recommendation}: keep US citizenship (the exit tax and loss of US access outweigh savings); plan the London move to use the four-year relief, realising US gains before arrival only if UK rules make that useful; set a review before year 9 in the UK to decide on leaving before inheritance tax exposure; and consider Singapore later for residence, not citizenship, coordinating with US filing.
\end{itemize}

\qffield{Conditions} Statutes and guidance as published by 2026. The four-year regime is described from the government's August 2024 policy summary; the enacted legislation governs the details. Residence rules have many exceptions; investment-migration programmes change often.

\qffield{Pitfalls} Assuming that moving abroad ends US tax for a citizen. Counting days without checking ties and treaty tie-breakers. Overlooking the inheritance-tax tail after leaving the UK.

\qffield{See also} \hyperref[meth:retirement-across-jurisdictions]{cross-border pensions}; \hyperref[meth:wealth-transfer-strategies]{transfers}.

\end{qfmethod}

\begin{qfmethod}[Complex families, business owners and executives: blended and cross-border families, entrepreneurs' concentration and liquidity events, qualified small business stock, executives' equity awards and trading plans]\label{meth:complex-client-situations}
\qffield{Statement} \emph{Complex families} (Interpretation): blended families (children from earlier marriages, the need to provide for a spouse and still protect children's inheritances); a non-citizen spouse (the US annual exclusion for gifts to such a spouse is USD~194,000 in 2026, not unlimited, \citep[sec.~4.42(2)]{irs-rev-proc-2025-32}); dependants with special needs; divorce and pre-marital agreements; and family members in several countries. Approaches: written wills and trusts that separate income for a spouse from capital for children, guardianship and support trusts, and early agreement on family rules (\hyperref[meth:advising-families]{family framework}).

\emph{Entrepreneurs and business owners}: most wealth sits in one illiquid, undiversified company whose fortunes are tied to the owner's human capital. Plans cover key-person and succession risk, separating personal from business assets, building liquid wealth outside the business, preparing for a sale or listing, and using tax provisions. In the US, gains on \emph{qualified small business stock} acquired after 4 July 2025 are excluded at 50, 75 or 100 percent after three, four or five years, up to the greater of USD~15 million and ten times basis per issuer \citep[sec.~1202(a)(1), (a)(5), (b)(1), (b)(4)]{us-internal-revenue-code-26}.

\emph{Professionals and executives}: concentrated employer stock and options, vesting schedules, insider-trading restrictions and liability exposure. A Rule~10b5-1 trading plan gives an affirmative defence if adopted in good faith without material nonpublic information; for directors and officers trading may start only after the later of 90 days and two business days after the next 10-Q or 10-K, but no more than 120 days after adoption, and they must certify the absence of inside information \citep[para.~(c)(1)(ii)(B)]{sec-rule-10b5-1}.

\qffield{Practice} (original; \texttt{private-wealth-checks.py}, part~12).
\begin{itemize}
\item \emph{Luis Rivera, founder}: holds shares in his company acquired after 4 July 2025 with a basis of 1 million. On a sale producing a 20 million gain after five years, 15 million is excluded (the greater of 15 million and 10 times basis), and the remaining 5 million is taxed at 23.8 percent: 1.19 million, against 4.76 million without the exclusion. Sold after three years, only 7.5 million would be excluded. \emph{Recommendation}: unless a buyer's offer is far better earlier, time the sale for the five-year mark; before it, build a liquid reserve outside the company and agree a succession plan through the family council.
\item \emph{An executive}: Tunde Okafor, if an officer, adopts a 10b5-1 plan on 4 March 2026. With the 10-Q filed 7 May, sales may start on 2 June; if the 10-Q were late (5 June), on 9 June; never after 2 July. The plan sells a fixed number of shares each month for 12 months.
\item \emph{A blended family}: Margaret Hartley's son David has remarried; his wife has two children, and he wants her provided for while his own assets eventually pass to his children. Recommendation: a life-interest trust for his wife with the capital passing to his children, a letter of wishes, and coordinated wills. Assets left to her outright would be exempt from inheritance tax but could pass to her children; the trust's tax treatment needs UK legal advice.
\end{itemize}

\qffield{Conditions} US federal law as amended in 2025 and SEC rules as in force in 2026; blended-family solutions depend on succession law in each jurisdiction.

\qffield{Pitfalls} Selling founder shares a few months before an exclusion improves. Adopting a trading plan while aware of inside information. Leaving everything outright to a second spouse and disinheriting children by accident.

\qffield{See also} \hyperref[meth:human-capital-risks]{human capital}; \hyperref[rule:standard-market-integrity]{Standard II}.

\end{qfmethod}

\begin{qfmethod}[Transferring wealth: lifetime gifts against bequests, US gift and estate tax and basis rules, UK inheritance tax and potentially exempt transfers, and philanthropy through trusts, donor funds and foundations]\label{meth:wealth-transfer-strategies}
\qffield{Statement} \emph{United States} (2026): gifts up to USD~19,000 per donee per year are excluded; above that, gifts use a unified lifetime exclusion that also shelters the estate, USD~15 million in 2026 and indexed afterwards, with the same amount exempt from generation-skipping tax \citep[secs.~3.14, 4.42]{irs-rev-proc-2025-32}. Estate and gift tax reaches 40 percent at the margin \citep[sec.~2001(c)]{us-internal-revenue-code-26}. Assets inherited take a basis equal to market value at death; gifted assets keep the donor's basis \citep[secs.~1014(a), 1015(a)]{us-internal-revenue-code-26}.

\emph{United Kingdom} (2026--27): inheritance tax at 40 percent above a GBP~325,000 nil-rate band, with an additional 175,000 when a home passes to direct descendants in estates below 2 million; 36 percent if 10 percent of the net estate goes to charity; unused bands pass to a surviving spouse; transfers to a spouse or charity are exempt \citep[overview]{govuk-iht-2026}; \citep[thresholds table]{govuk-iht-thresholds-2026}. Gifts are exempt if the donor survives seven years; on death within seven years, tax on gifts above the nil-rate band tapers from 40 percent (under 3 years) to 32, 24, 16 and 8 percent; an annual exemption of 3,000, small gifts of 250 and regular gifts from income are exempt; gifts with reservation of benefit stay in the estate \citep[gifts]{govuk-iht-2026}.

\emph{Philanthropy across generations}. US options include outright gifts, donor-advised funds, private foundations, which must distribute about 5 percent of their non-charitable assets a year \citep[sec.~4942(e)]{us-internal-revenue-code-26}, and charitable remainder trusts, which pay the donor a fixed 5--50 percent annuity (or a unitrust amount) for life or up to 20 years and leave the remainder to charity, the remainder being worth at least 10 percent at creation \citep[sec.~664(d)]{us-internal-revenue-code-26}. Foundations also serve as a vehicle for family governance and education (Interpretation).

\qffield{Practice} (original; \texttt{private-wealth-checks.py}, part~13).
\begin{itemize}
\item \emph{Gift or bequest of a growth asset} (US). An asset worth 5 million with a 1 million basis is expected to be worth 10 million at death; the estate has 30 million of other assets. Kept, it is taxed in the estate: heirs receive 30.0 million after 10.0 million of estate tax, with the asset's basis stepped up. Given now, it uses 5 million of the exemption, estate tax falls to 8.0 million and heirs receive 32.0 million if they keep the asset, but 29.86 million if they sell and pay 23.8 percent on the 9 million carryover gain. \emph{Decision}: give it now only if the family will hold it or it will grow faster; in an estate below the exemption, keep it and let the basis step up.
\item \emph{Annual exclusions}: 19,000 to each of six grandchildren for ten years moves 1.14 million out of the estate without using exemption.
\item \emph{UK gift} (Margaret Hartley): 500,000 to her son, who survives her 4.5 years: the first 325,000 uses her nil-rate band; the remaining 175,000 is taxed at 24 percent, 42,000, against 70,000 if she died within three years.
\item \emph{UK estate}: her 2.4 million estate including the home is above 2 million, so only the 325,000 band applies: 830,000 of tax at 40 percent. Leaving 10 percent of the net estate to charity lowers the rate to 36 percent; lifetime gifts survived by seven years reduce the estate.
\item \emph{Charitable trust} (US): a 1 million charitable remainder annuity trust paying 50,000 a year for 20 years: the payments are worth 623,000 at 5 percent, leaving a remainder of 377,000, above the 10 percent minimum.
\end{itemize}

\qffield{Conditions} The US estate figure assumes a flat 40 percent above the exemption and no state tax; the remainder valuation uses a 5 percent rate rather than the IRS rate for the month. UK rules as on GOV.UK in 2026.

\qffield{Pitfalls} Gifting low-basis assets to heirs who will sell them. Keeping a benefit in a gifted home. Treating charitable trusts as tax shelters rather than as gifts.

\qffield{See also} \hyperref[meth:residency-citizenship]{residence and inheritance tax}; \hyperref[meth:advising-families]{family governance}.

\end{qfmethod}
